% Nota de noc_noc. GERADA por flim_al/tabelas.py.
\footnotesize NoC@X é o número de cliques até a IoU passar do alvo, com teto de 12; é o eixo de RITM (arXiv:2102.06583) e SimpleClick (arXiv:2210.11006), e \textbf{não} se compara com o NoC publicado deles, que usa outros conjuntos e outra definição de alvo. Imagem que não atinge o alvo entra com o teto, e por isso a \textbf{taxa de sucesso} aparece em toda linha: um braço pode mostrar NoC baixo por desistir mais cedo. \texttt{$\alpha$=1} é o FLIM puro, que refaz o banco de filtros a cada clique; \texttt{$\alpha$=0} congela o banco, e o clique chega ao modelo só pelos rótulos que o decoder \texttt{labeled\_marker} lê. Os testes são \textbf{separados por dataset}: agregar Schisto e BraTS, com 7\% e 30\% de sucesso, inflou um achado exploratório a p=0,0265 que separado não passa de p=0,10. O clique é simulado do ground truth pelo protocolo de Xu et al. (2016), com usuário que acerta sempre, então a tabela mede \textbf{número de interações e não tempo de especialista}.
\normalsize
